% ============================================================================
% LANGENTWURF LE-020: Probleme lösen
% Profil: S (Senioren 65+) | Tier: 2 (Standard)
% ============================================================================

\documentclass[11pt,a4paper]{article}
\usepackage[utf8]{inputenc}
\usepackage[T1]{fontenc}
\usepackage[ngerman]{babel}
\usepackage{geometry}
\usepackage{fancyhdr}
\usepackage{lastpage}
\usepackage{xcolor}
\usepackage{tcolorbox}
\usepackage{enumitem}
\usepackage{longtable}
\usepackage{hyperref}
\usepackage{amssymb}

\geometry{left=2.5cm, right=2.5cm, top=2.5cm, bottom=2.5cm}
\definecolor{fach}{HTML}{b35c00}
\definecolor{gray}{HTML}{666666}
\definecolor{successgreen}{HTML}{2a7a4b}
\definecolor{warnred}{HTML}{c62828}

\tcbuselibrary{skins,breakable}
\newtcolorbox{scorebox}{ colback=white, colframe=fach, fonttitle=\bfseries, title=Didaktik-Score, breakable }
\newtcolorbox{tippbox}[1][]{ colback=successgreen!5, colframe=successgreen, fonttitle=\bfseries, title=#1, breakable }

\pagestyle{fancy}
\fancyhf{}
\fancyhead[L]{\small\textcolor{gray}{Digitale Grundlagen -- 020 Probleme lösen}}
\fancyhead[R]{\small\textcolor{gray}{LE Tier 2 / Profil S}}
\fancyfoot[C]{\small\thepage{} / \pageref{LastPage}}
\renewcommand{\headrulewidth}{0.4pt}

\begin{document}

\begin{titlepage}
    \centering
    \vspace*{2cm}
    {\Large\textcolor{gray}{Digitale Grundlagen -- Senioren 65+}}\\[2cm]
    {\Huge\textcolor{fach}{\textbf{Langentwurf}}}\\[0.5cm]
    {\large Tier 2 | Profil S}\\[2cm]
    {\LARGE\textbf{020 -- Probleme selbst lösen}}\\[0.5cm]
    {\large Modul E: Hilfe \& Selbstständigkeit}\\[2cm]
    \begin{tabular}{rl}
        \textbf{Zeitrahmen:} & 60--75 Minuten \\
        \textbf{Vorkenntnisse:} & Grundbedienung iPhone/iPad \\
    \end{tabular}
    \vfill
    {\normalsize Stand: \today}
\end{titlepage}

\tableofcontents
\newpage

\section{Bedingungsanalyse}

\subsection{Ausgangssituation}
\begin{itemize}
    \item Panik bei jeder Fehlermeldung
    \item ``Es geht nicht'' ohne genaue Beschreibung
    \item Neustart als Lösung unbekannt
    \item Angst, etwas kaputt zu machen
    \item Bei jedem Problem sofort Hilfe rufen
\end{itemize}

\subsection{Typische Probleme}
\begin{itemize}
    \item ``Mein Handy ist so langsam''
    \item ``Der Speicher ist voll''
    \item ``Das Internet geht nicht''
    \item ``Eine App reagiert nicht mehr''
\end{itemize}

\section{Sachanalyse}

\subsection{Kernkonzepte}
\begin{enumerate}
    \item \textbf{Neustart:} Löst 80\% aller Probleme (``Aus- und wieder einschalten'')
    \item \textbf{App schließen:} Hängende App beenden
    \item \textbf{Speicher prüfen:} Warum ist das Gerät voll?
    \item \textbf{WLAN prüfen:} Ist überhaupt Internet da?
    \item \textbf{Updates:} Manchmal fehlt nur ein Update
\end{enumerate}

\subsection{Alltagsanalogien}
\begin{tabular}{|l|p{8cm}|}
\hline
\textbf{Begriff} & \textbf{Analogie} \\
\hline
Neustart & Kurz schlafen und frisch aufwachen -- alles neu sortiert \\
App schließen & Tür zumachen und wieder öffnen \\
Speicher voll & Schrank vollgestopft -- erst aufräumen \\
WLAN-Verbindung & Telefonleitung -- ist sie überhaupt da? \\
Update & Verbesserungen vom Hersteller -- wie Wartung beim Auto \\
\hline
\end{tabular}

\section{Didaktische Analyse}

\subsection{Stellung in der Reihe}
Erste Einheit im Modul ``Hilfe \& Selbstständigkeit''. Ermutigt zur Eigenständigkeit.

\subsection{Didaktische Reduktion}
\textbf{Ausgeklammert:} Tiefere Fehlersuche, Systemeinstellungen, Recovery-Modus.

\textbf{Fokus:} Die 5 einfachsten Selbsthilfe-Schritte, die 90\% der Alltagsprobleme lösen.

\section{Lernziele}

\begin{tabular}{|c|p{7cm}|c|c|}
\hline
\textbf{Nr.} & \textbf{Die Teilnehmer können...} & \textbf{Stufe} & \textbf{Niveau} \\
\hline
FZ1 & das Gerät neu starten & Anwenden & Basis \\
FZ2 & eine hängende App schließen & Anwenden & Basis \\
FZ3 & den Speicherstand prüfen & Anwenden & Basis \\
FZ4 & die WLAN-Verbindung prüfen & Anwenden & Basis \\
FZ5 & nach Updates suchen & Anwenden & Vertiefung \\
\hline
\end{tabular}

\section{Methodische Analyse}

\subsection{Die 5 Erste-Hilfe-Schritte}
\begin{tippbox}[Wenn etwas nicht funktioniert]
\begin{enumerate}
    \item \textbf{Ruhe bewahren} -- Du machst nichts kaputt!
    \item \textbf{Neustart} -- Gerät aus- und wieder einschalten
    \item \textbf{App schließen} -- Wenn nur eine App hängt
    \item \textbf{WLAN prüfen} -- Ist Internet da?
    \item \textbf{Speicher prüfen} -- Ist noch Platz?
\end{enumerate}
\end{tippbox}

\section{Verlaufsplanung}

\begin{longtable}{|p{1cm}|p{2cm}|p{5cm}|p{3cm}|p{2cm}|}
\hline
\textbf{Zeit} & \textbf{Phase} & \textbf{Inhalt} & \textbf{TN-Aktivität} & \textbf{Material} \\
\hline
0--10 & Einstieg & ``Was tun Sie, wenn Ihr Fernseher nicht geht?'' Parallele zu Technik & Berichten & -- \\
\hline
10--25 & Neustart & iPhone/iPad neu starten demonstrieren und üben & Selbst üben & S.1 \\
\hline
25--40 & App schließen & App-Umschalter zeigen, App ``wegwischen'' & Selbst üben & S.1 \\
\hline
40--55 & Prüfen & WLAN und Speicher prüfen, Einstellungen finden & Selbst prüfen & S.2 \\
\hline
55--70 & Updates & Nach Updates suchen, Auto-Updates erklären & Anschauen & S.2-3 \\
\hline
\end{longtable}

\section{Lösungen}

\textbf{iPhone/iPad neu starten:}
\begin{itemize}
    \item \textbf{Mit Face ID:} Seitentaste + Lautstärke-Taste drücken, ``Ausschalten'' wischen, dann Seitentaste zum Einschalten
    \item \textbf{Mit Home-Button:} Seitentaste lange drücken, ``Ausschalten'' wischen, dann wieder einschalten
    \item Warten, bis der Apfel erscheint!
\end{itemize}

\textbf{App schließen:}
\begin{itemize}
    \item Von unten nach oben wischen und in der Mitte stoppen (App-Umschalter)
    \item Die hängende App nach oben ``wegwischen''
    \item App neu öffnen
\end{itemize}

\textbf{WLAN prüfen:}
\begin{itemize}
    \item Einstellungen $\rightarrow$ WLAN
    \item Ist WLAN an (grüner Schalter)?
    \item Ist ein Netzwerk verbunden (Häkchen)?
    \item Tipp: WLAN aus- und wieder einschalten
\end{itemize}

\textbf{Speicher prüfen:}
\begin{itemize}
    \item Einstellungen $\rightarrow$ Allgemein $\rightarrow$ iPhone-Speicher
    \item Farbiger Balken zeigt Belegung
    \item Was verbraucht am meisten? (Fotos? Apps?)
\end{itemize}

\textbf{Updates prüfen:}
\begin{itemize}
    \item Einstellungen $\rightarrow$ Allgemein $\rightarrow$ Softwareupdate
    \item Wenn Update verfügbar: ``Laden und installieren''
    \item Tipp: ``Automatische Updates'' aktivieren
\end{itemize}

\section{Didaktik-Score}

\begin{scorebox}
\begin{tabular}{|c|l|c|c|p{3.5cm}|}
\hline
\# & Dimension & Gew. & Score & Notiz \\
\hline
1 & Differenzierung & 15\% & 8/10 & Basis/Vertiefung klar \\
2 & Taxonomie & 10\% & 9/10 & Stark Anwenden \\
3 & Alltagsrelevanz & 10\% & 10/10 & Täglich gebraucht! \\
4 & Aktivierung & 10\% & 10/10 & Sofort selbst üben \\
5 & Scaffolding & 10\% & 10/10 & 5-Schritte-Anleitung \\
6 & Feedback & 10\% & 9/10 & Ermutigung stark \\
7 & Sprachliche Klarheit & 10\% & 9/10 & Einfache Sprache \\
8 & Motivation & 10\% & 10/10 & Selbstständigkeit! \\
9 & Barrierefreiheit & 10\% & 9/10 & Große Darstellung \\
10 & Praktikabilität & 5\% & 9/10 & 70 Min gut \\
\hline
\multicolumn{3}{|r|}{\textbf{GESAMT:}} & \textbf{9.35/10} & \textcolor{successgreen}{\textbf{FREIGABE}} \\
\hline
\end{tabular}
\end{scorebox}

\end{document}
